% GENERATED FILE. Edit canonical lessons, then run npm run generate:books.
% canonical-source-hash: fnv1a64:8184a46dedbd1469
% canonical-lessons: LA-C351-concido, LA-C351-seco, LA-C351-mortarium, LA-C351-pistillum, LA-C351-tero

\chapter{Concīdere et terere}
\label{ch:la-concidere-et-terere}
\hlchaptermodality{351}

\begin{chapteropening}
\textbf{By the end of this chapter:} \emph{I can describe preparing food: to chop up, to cut, the mortar, the pestle and to grind.}

Complete the last lesson of chapter 351: I can describe preparing food: to chop up, to cut, the mortar, the pestle and to grind.
\end{chapteropening}

\section[concīdō, concīdere]{concīdō, concīdere — to chop up, to cut to pieces}

\label{lesson:LA-C351-concido}



\begin{quote}
Before the new one: say the Latin for to wipe, then the Latin for a small knife.
\end{quote}

\subsection*{You'll want to know: concīdō, concīdere}
\textbf{concīdō, concīdere} — \textquotedblleft{}to chop up, to cut to pieces\textquotedblright{}.

\begin{quote}
Coquus cēpās et ālium concīdit.
\end{quote}

\begin{quote}
\emph{The cook chopped up the onions and garlic.}
\end{quote}

\emph{concīdō, concīdere, concīdī, concīsum}, long \emph{ī}: \emph{con-} + \emph{caedō}. Do not confuse \emph{concidō}, 'to collapse'. English \emph{concise}.

\begin{cousinweb}
\textbf{concīdō} is built on \emph{caedō}, \textquotedblleft{}to cut\textquotedblright{}, with \emph{con-} (EtymDB): to cut up small. English \emph{concise}, cut short, comes from its participle.
\end{cousinweb}

\begin{grammarlens}[title={What you've built}]
One more word for telling, explaining and describing at B1.
\end{grammarlens}

\subsection*{Your turn}
\begin{itemize}
\raggedright
  \item \emph{Say it:} \emph{concīdō, concīdere}
  \item \emph{Say it:} \emph{concīdō, concīdere}, once more
  \item \emph{Say it:} a sentence of your own with \emph{concīdō, concīdere}
  \item \emph{Recall:} say the Latin for to wipe, then the Latin for a small knife, then say \emph{concīdō, concīdere} again
\end{itemize}

\begin{tcolorbox}[breakable,colback=teal!4,colframe=teal!35!black,title={Before you move on}]
What does concīdō, concīdere mean? (\textquotedblleft{}To chop up, to cut to pieces\textquotedblright{}.) Say it once more.
\end{tcolorbox}
\section[secō, secāre]{secō, secāre — to cut, to slice}

\label{lesson:LA-C351-seco}



\begin{quote}
Before the new one: say the Latin for a small knife, then the Latin for to chop up.
\end{quote}

\subsection*{You'll want to know: secō, secāre}
\textbf{secō, secāre} — \textquotedblleft{}to cut, to slice\textquotedblright{}.

\begin{quote}
Pater pānem cultellō secuit.
\end{quote}

\begin{quote}
\emph{Father sliced the bread with a little knife.}
\end{quote}

\emph{secō, secāre, secuī, sectum}. English \emph{section} and \emph{insect}, a creature cut into parts, come from it.

\begin{cousinweb}
\textbf{secō} is inherited from the Proto-Indo-European root \emph{*sek-}, \textquotedblleft{}to cut\textquotedblright{} (EtymDB). English \emph{section}, \emph{sector} and \emph{insect} come from it.
\end{cousinweb}

\begin{grammarlens}[title={What you've built}]
One more word for telling, explaining and describing at B1.
\end{grammarlens}

\subsection*{Your turn}
\begin{itemize}
\raggedright
  \item \emph{Say it:} \emph{secō, secāre}
  \item \emph{Say it:} \emph{secō, secāre}, once more
  \item \emph{Say it:} a sentence of your own with \emph{secō, secāre}
  \item \emph{Recall:} say the Latin for a small knife, then the Latin for to chop up, then say \emph{secō, secāre} again
\end{itemize}

\begin{tcolorbox}[breakable,colback=teal!4,colframe=teal!35!black,title={Before you move on}]
What does secō, secāre mean? (\textquotedblleft{}To cut, to slice\textquotedblright{}.) Say it once more.
\end{tcolorbox}
\section[mortārium, mortāriī]{mortārium, mortāriī — a mortar (a bowl for pounding)}

\label{lesson:LA-C351-mortarium}



\begin{quote}
Before the new one: say the Latin for to chop up, then the Latin for to cut.
\end{quote}

\subsection*{You'll want to know: mortārium, mortāriī}
\textbf{mortārium, mortāriī} — \textquotedblleft{}a mortar (a bowl for pounding)\textquotedblright{}.

\begin{quote}
Coquus herbās in mortārium posuit.
\end{quote}

\begin{quote}
\emph{The cook put the herbs into the mortar.}
\end{quote}

Neuter. A heavy bowl of stone or clay. English \emph{mortar}, for building too, comes from it.

\begin{cousinweb}
\textbf{mortārium} is a Latin word whose origin is not settled (Lewis \& Short). English \emph{mortar}, both the bowl and the builder mix, comes from it.
\end{cousinweb}

\begin{grammarlens}[title={What you've built}]
One more word for telling, explaining and describing at B1.
\end{grammarlens}

\subsection*{Your turn}
\begin{itemize}
\raggedright
  \item \emph{Say it:} \emph{mortārium, mortāriī}
  \item \emph{Say it:} \emph{mortārium, mortāriī}, once more
  \item \emph{Say it:} a sentence of your own with \emph{mortārium, mortāriī}
  \item \emph{Recall:} say the Latin for to chop up, then the Latin for to cut, then say \emph{mortārium, mortāriī} again
\end{itemize}

\begin{tcolorbox}[breakable,colback=teal!4,colframe=teal!35!black,title={Before you move on}]
What does mortārium, mortāriī mean? (\textquotedblleft{}A mortar (a bowl for pounding)\textquotedblright{}.) Say it once more.
\end{tcolorbox}
\section[pistillum, pistillī]{pistillum, pistillī — a pestle}

\label{lesson:LA-C351-pistillum}



\begin{quote}
Before the new one: say the Latin for to cut, then the Latin for a mortar.
\end{quote}

\subsection*{You'll want to know: pistillum, pistillī}
\textbf{pistillum, pistillī} — \textquotedblleft{}a pestle\textquotedblright{}.

\begin{quote}
Pistillum ligneum erat et grave.
\end{quote}

\begin{quote}
\emph{The pestle was wooden and heavy.}
\end{quote}

Neuter, from a verb meaning 'to pound'. English \emph{pestle} comes from it.

\begin{cousinweb}
\textbf{pistillum} is built on \emph{pīnsō}, \textquotedblleft{}to pound\textquotedblright{} (EtymDB), with a small-tool ending. English \emph{pestle} comes from it through French, and \emph{pistil} directly.
\end{cousinweb}

\begin{grammarlens}[title={What you've built}]
One more word for telling, explaining and describing at B1.
\end{grammarlens}

\subsection*{Your turn}
\begin{itemize}
\raggedright
  \item \emph{Say it:} \emph{pistillum, pistillī}
  \item \emph{Say it:} \emph{pistillum, pistillī}, once more
  \item \emph{Say it:} a sentence of your own with \emph{pistillum, pistillī}
  \item \emph{Recall:} say the Latin for to cut, then the Latin for a mortar, then say \emph{pistillum, pistillī} again
\end{itemize}

\begin{tcolorbox}[breakable,colback=teal!4,colframe=teal!35!black,title={Before you move on}]
What does pistillum, pistillī mean? (\textquotedblleft{}A pestle\textquotedblright{}.) Say it once more.
\end{tcolorbox}
\section[terō, terere]{terō, terere — to grind, to rub}

\label{lesson:LA-C351-tero}



\begin{quote}
Before the new one: say the Latin for a mortar, then the Latin for a pestle.
\end{quote}

\subsection*{You'll want to know: terō, terere}
\textbf{terō, terere} — \textquotedblleft{}to grind, to rub\textquotedblright{}.

\begin{quote}
Ancilla piper pistillō terēbat.
\end{quote}

\begin{quote}
\emph{The maid was grinding pepper with the pestle.}
\end{quote}

\emph{terō, terere, trīvī, trītum}. English \emph{trite}, worn out by rubbing.

\begin{cousinweb}
\textbf{terō} is an inherited verb; Lewis \& Short compare Greek \emph{teirō}, \textquotedblleft{}to wear away\textquotedblright{}. English \emph{trite} comes from its participle \emph{trītus}.
\end{cousinweb}

\begin{grammarlens}[title={What you've built}]
One more word for telling, explaining and describing at B1.
\end{grammarlens}

\subsection*{Your turn}
\begin{itemize}
\raggedright
  \item \emph{Say it:} \emph{terō, terere}
  \item \emph{Say it:} \emph{terō, terere}, once more
  \item \emph{Say it:} a sentence of your own with \emph{terō, terere}
  \item \emph{Recall:} say the Latin for a mortar, then the Latin for a pestle, then say \emph{terō, terere} again
\end{itemize}

\begin{tcolorbox}[breakable,colback=teal!4,colframe=teal!35!black,title={Before you move on}]
What does terō, terere mean? (\textquotedblleft{}To grind, to rub\textquotedblright{}.) Say it once more.
\end{tcolorbox}
